% !TEX root = main.tex
\section{Phát biểu bài toán}\label{sec:problem}

\subsection{Mô tả bài toán}\label{sec:problem:desc}

Cho mạng đường lái xe của quận Hà Đông, một điểm đi, một điểm đến và một giờ xuất phát. Vận tốc trên mỗi đoạn đường
thay đổi theo giờ trong ngày theo một quy luật lặp lại; ta có trong tay một \emph{dự báo} của quy luật đó (gọi là vận
tốc nền $V_{\text{base}}$). Tuy nhiên ngày thực tế luôn lệch khỏi dự báo: có nhiễu (đường đông hơn hoặc vắng hơn bình
thường, và độ lệch này lan sang các đoạn lân cận, kéo dài sang các giờ kế tiếp) và có thể xảy ra \emph{sự cố} làm một
đoạn đường chậm hẳn trong 30--120 phút. Người lái chỉ \emph{quan sát} được vận tốc thực trong một bán kính quanh vị trí
hiện tại; trong một số thiết lập, có thêm dữ liệu từ các cảm biến cố định trên một phần mạng đường.

Nhiệm vụ là chọn tuyến đường --- hoặc chọn một lần lúc xuất phát, hoặc điều chỉnh liên tục trong khi đi --- sao cho đến
nơi sớm nhất. Bên cạnh đó có ba bài toán phụ: (i) có nên đợi thêm một khoảng thời gian trước khi xuất phát để tránh
cao điểm; (ii) chọn tuyến \emph{tin cậy} (ít rủi ro bị trễ) thay vì tuyến nhanh nhất trung bình; (iii) khi nhiều xe
cùng được dẫn đường, chỉ dẫn nên được phân bổ thế nào để không tự tạo ra ùn tắc mới.

\textbf{Đầu vào} của một truy vấn gồm: nút đi $s$, nút đến $d$, giờ xuất phát $t_0$; đầu vào của hệ thống gồm đồ thị
đường, vận tốc nền 24 khung giờ, và (tuỳ phương pháp) các quan sát thời gian thực. \textbf{Đầu ra} là dãy cạnh từ $s$ tới
$d$ (hoặc một chính sách chọn cạnh tiếp theo), cùng thời gian di chuyển thực tế phát sinh khi đi theo nó.

\begin{vd}[Tuyến ngắn nhất không phải lúc nào cũng nhanh nhất]\label{vd:toy}
Xét một ví dụ \emph{giả định} để minh hoạ: từ $s$ tới $d$ có hai lựa chọn. Tuyến 1 đi trục chính dài $3{,}0$~km; tuyến 2
đi đường nhánh dài $2{,}4$~km. Lúc 3 giờ sáng, vận tốc trục chính là $50$~km/h và đường nhánh $18$~km/h, nên
\begin{equation}
T_1(3\text{h}) = \frac{3{,}0}{50}\cdot 60 = 3{,}6 \text{ phút}, \qquad
T_2(3\text{h}) = \frac{2{,}4}{18}\cdot 60 \approx 8{,}0 \text{ phút},
\label{eq:toy3}
\end{equation}
tuyến 1 (dài hơn) nhanh hơn hẳn. Lúc 8 giờ sáng, trục chính tắc còn $14$~km/h, đường nhánh còn $13$~km/h:
\begin{equation}
T_1(8\text{h}) = \frac{3{,}0}{14}\cdot 60 \approx 12{,}9 \text{ phút}, \qquad
T_2(8\text{h}) = \frac{2{,}4}{13}\cdot 60 \approx 11{,}1 \text{ phút},
\label{eq:toy8}
\end{equation}
tuyến 2 trở thành tuyến nhanh nhất. Tuyến tối ưu phụ thuộc vào giờ xuất phát.
\end{vd}

\begin{vd}[Ví dụ thật trong tập kiểm tra]\label{vd:real}
Truy vấn \ExampleQuery{} trong tập kiểm tra có \ExampleNInc{} sự cố trên hành lang đi (Hình~\ref{fig:route} ở
Phần~\ref{sec:exp}). Kế hoạch tốt nhất theo dự báo (plan-once) đi thẳng vào đoạn bị sự cố và mất \ExamplePoTT{} phút;
cơ chế tái lập kế hoạch nhận ra đoạn chậm khi đến gần và vòng tránh, chỉ mất \ExampleRpTT{} phút --- gần bằng lời giải
của một ``nhà tiên tri'' biết trước toàn bộ tương lai (oracle, \ExampleOrTT{} phút). Ví dụ cho thấy ba yếu tố cốt lõi
của bài toán: phụ thuộc thời gian, bất định, và quan sát cục bộ.
\end{vd}

\subsection{Mô hình hoá toán học}\label{sec:problem:model}

\textbf{Đồ thị và hàm thời gian.} Mạng đường được biểu diễn bởi đồ thị có hướng $G=(V,E)$ với $|V|=n$ nút (giao lộ,
điểm cuối đường) và $|E|=m$ cạnh (đoạn đường một chiều). Mỗi cạnh $e=(u,v)\in E$ có chiều dài $L_e$ (mét). Thời gian
$t\in\mathbb{R}$ tính bằng phút kể từ 0:00; vận tốc lặp lại với chu kỳ $1440$ phút. Ký hiệu $v_e(t)>0$ (m/phút) là vận tốc
trên cạnh $e$ tại thời điểm $t$.

\begin{dn}[Hàm thời gian qua cạnh và hàm đến]\label{dn:arrival}
Thời gian để đi hết cạnh $e$ khi vào cạnh tại thời điểm $t$ và thời điểm ra khỏi cạnh lần lượt là
\begin{equation}
\tau_e(t) = \frac{L_e}{v_e(t)},
\label{eq:tau}
\end{equation}
\begin{equation}
a_e(t) = t + \tau_e(t).
\label{eq:arrival}
\end{equation}
\end{dn}

Một đường đi từ $s$ tới $d$ là dãy cạnh $P=(e_1,\dots,e_k)$ với $e_i=(u_{i-1},u_i)$, $u_0=s$, $u_k=d$. Thời điểm đến
khi xuất phát lúc $t_0$ được tính bằng cách hợp các hàm đến:
\begin{equation}
A(P,t_0) = a_{e_k}\circ a_{e_{k-1}}\circ\cdots\circ a_{e_1}(t_0).
\label{eq:path-arrival}
\end{equation}

\begin{dn}[Tính FIFO]\label{dn:fifo}
Mạng thoả tính \emph{FIFO} (first-in-first-out, đi trước thì đến trước) nếu với mọi cạnh $e$ và mọi $t_1\le t_2$:
\begin{equation}
a_e(t_1)\le a_e(t_2).
\label{eq:fifo}
\end{equation}
Khi $\tau_e$ khả vi, điều kiện \eqref{eq:fifo} tương đương với
\begin{equation}
\frac{d\tau_e(t)}{dt} \ge -1 \qquad \forall t.
\label{eq:fifo-deriv}
\end{equation}
\end{dn}

\textbf{Bài toán TDSP (lập kế hoạch một lần).} Với tập $\mathcal{P}(s,d)$ các đường đi từ $s$ tới $d$:
\begin{equation}
P^\ast = \arg\min_{P\in\mathcal{P}(s,d)} \bigl[A(P,t_0)-t_0\bigr],
\label{eq:tdsp}
\end{equation}
trong đó $v_e(t)$ lấy từ dự báo $V_{\text{base}}$ (phương pháp \emph{plan-once}) hoặc từ môi trường thực $V_{\text{real}}$
(lời giải \emph{oracle}, chỉ dùng làm cận dưới để đánh giá). Ràng buộc duy nhất là $P$ phải là đường đi hợp lệ trên $G$
(các cạnh nối tiếp nhau); không cho phép chờ tại nút.

\textbf{Bài toán trực tuyến (tái lập kế hoạch, học tăng cường).} Trạng thái đầy đủ tại bước $k$ gồm vị trí $u_k$, thời
điểm $t_k$ và toàn bộ trường vận tốc thực $V_{\text{real}}$; nhưng tác tử chỉ quan sát được
\begin{equation}
o_k = \Bigl(u_k,\; t_k,\; \bigl\{\tau^{\text{real}}_e(t_k) : e\in\mathcal{N}_r(u_k)\bigr\}\Bigr),
\label{eq:observation}
\end{equation}
với $\mathcal{N}_r(u)$ là tập cạnh có nút đầu cách $u$ không quá $r$ mét. Bài toán là một quá trình quyết định Markov
quan sát một phần (POMDP): tìm chính sách $\pi$ ánh xạ lịch sử quan sát thành cạnh ra tiếp theo để cực tiểu kỳ vọng thời
gian đi
\begin{equation}
\pi^\ast = \arg\min_\pi \; \mathbb{E}\bigl[T^\pi\bigr], \qquad T^\pi = A(P^\pi,t_0)-t_0 .
\label{eq:pomdp}
\end{equation}
Chất lượng của một phương pháp trên một truy vấn được đo bằng \emph{regret} so với oracle:
\begin{equation}
\text{Regret}(\pi) = 100\cdot\left(\frac{T^\pi}{T^{\text{oracle}}}-1\right)\ (\%).
\label{eq:regret}
\end{equation}

\textbf{Biến thể theo rủi ro.} Khi $V_{\text{real}}$ là biến ngẫu nhiên với phân phối $\mathcal{D}$, ta có thể ưu tiên
tuyến \emph{tin cậy} bằng cách cực tiểu phân vị 90\% thay vì kỳ vọng:
\begin{equation}
P^{\ast}_{0.9} = \arg\min_{P\in\mathcal{P}(s,d)} Q_{0.9}\bigl[A(P,t_0)-t_0\bigr], \qquad V_{\text{real}}\sim\mathcal{D}.
\label{eq:risk}
\end{equation}

\textbf{Chọn giờ xuất phát.} Với bước chờ $\Delta$ phút, tối đa $K$ lần chờ và trọng số chi phí chờ $w\in(0,1]$:
\begin{equation}
k^\ast = \arg\min_{k\in\{0,\dots,K\}} \Bigl[w\cdot k\Delta + T^{\text{oracle}}(t_0+k\Delta)\Bigr].
\label{eq:departure}
\end{equation}
Với $w=1$, đây đúng là ``tổng thời gian chờ và đi''; $w<1$ mô tả trường hợp chờ ở nhà/cơ quan ít tốn kém hơn ngồi trên xe.

\textbf{Nhiều xe.} Khi $N$ xe được dẫn đường cùng lúc, tải dẫn đường $q_e(t)$ (xe/giờ) làm chậm cạnh theo hàm BPR của
Cục Đường bộ Hoa Kỳ~\cite{bpr1964}:
\begin{equation}
\tau^{\text{load}}_e(t) = \tau_e(t)\left[1+\alpha\left(\frac{q_e(t)}{c_e}\right)^{\beta}\right],
\label{eq:bpr}
\end{equation}
với $c_e$ là năng lực thông hành của cạnh; mục tiêu hệ thống là cực tiểu tổng (hay trung bình) thời gian đi của các xe.

\textbf{Đặc tính của bài toán.}
\begin{dl}[\cite{dreyfus1969,kaufman1993}]\label{dl:poly}
Nếu mạng thoả FIFO \eqref{eq:fifo} và $\tau_e(t)>0$, bài toán \eqref{eq:tdsp} giải được trong thời gian đa thức bằng giải
thuật gắn nhãn kiểu Dijkstra theo thời điểm đến (chứng minh ở Phần~\ref{sec:method:correct}).
\end{dl}

\begin{nx}
Khi mất tính FIFO và không cho phép chờ tại nút, đường tối ưu có thể không là đường đơn và tính chất cấu trúc con tối
ưu bị phá vỡ~\cite{orda1990}. Vì vậy trong mô phỏng chúng tôi \emph{ép} tính FIFO (Phần~\ref{sec:method:speed}).
\end{nx}

\begin{nx}
Bài toán trực tuyến \eqref{eq:pomdp} là một POMDP với không gian trạng thái liên tục và rất lớn (trường vận tốc trên
$m$ cạnh); không có lời giải dạng đóng. Các phương pháp trong đề tài là heuristic có cơ sở (tái lập kế hoạch) hoặc học
từ dữ liệu (RL).
\end{nx}

\begin{nx}
Mục tiêu phân vị \eqref{eq:risk} \emph{không cộng tính} theo cạnh: phân vị của tổng khác tổng các phân vị, nên bài toán
mất tính chất cấu trúc con tối ưu và không giải trực tiếp được bằng Dijkstra. Đề tài dùng tập tuyến ứng viên
(Yen~\cite{yen1971}) rồi chọn theo phân vị ước lượng.
\end{nx}

\begin{nx}
Bài toán nhiều xe là một bài toán phân bổ giao thông: nếu mỗi xe tự chọn tuyến nhanh nhất cho mình, hệ thống hội tụ về
cân bằng người dùng (nguyên lý Wardrop~\cite{wardrop1952}), nói chung khác với tối ưu toàn hệ thống.
\end{nx}

\subsection{Khó khăn và thách thức}\label{sec:problem:challenges}

\textbf{Độ phức tạp.} Mạng đường Hà Đông sau làm sạch có $n=\NNodes$ nút và $m=\NEdges$ cạnh. Tìm kiếm phụ thuộc thời
gian một lần có độ phức tạp $O((n+m)\log n)$ --- rất nhanh; nhưng các bài toán xung quanh nhân chi phí đó lên nhiều lần:
tái lập kế hoạch chạy lại tìm kiếm tại mỗi giao lộ (trung bình hàng chục lần mỗi chuyến), định tuyến theo rủi ro đánh
giá nhiều ứng viên trên nhiều kịch bản, mô phỏng nhiều xe chạy hàng trăm tìm kiếm có tương tác, còn RL cần hàng trăm
nghìn đến hàng triệu bước tương tác để học.

\textbf{Các khía cạnh khiến bài toán khó.}
\begin{enumerate}[label=(\arabic*)]
    \item \emph{Không có dữ liệu vận tốc thực} cho Hà Đông: phải xây dựng mô hình mô phỏng có cấu trúc và minh bạch giả
    định; mọi kết luận định lượng chỉ có giá trị trong mô phỏng.
    \item \emph{Bảo đảm FIFO}: khi nội suy vận tốc theo phút, một đoạn đường dài có vận tốc tăng nhanh có thể làm
    $d\tau/dt<-1$; khi sự cố kết thúc, vận tốc nhảy vọt cũng phá vỡ FIFO.
    \item \emph{Bất định và quan sát cục bộ}: dự báo sai có hệ thống (sự cố) chỉ được phát hiện khi xe đã đến gần.
    \item \emph{Khó khăn riêng của RL}: phần thưởng thưa (chỉ biết đã đến đích ở cuối chuyến), tầm nhìn dài (một chuyến
    gồm hàng chục cạnh), \DeadEndPct\% số nút là ngõ cụt, tập hành động thay đổi theo nút (cần mặt nạ hành động), và chính
    sách phải tổng quát hoá sang cặp điểm đi--đến, giờ xuất phát và vị trí sự cố chưa từng thấy.
    \item \emph{Đánh giá công bằng}: các phương pháp phải chạy trên cùng tập kiểm tra, cùng môi trường thực tái lập theo
    seed, tách tập huấn luyện/kiểm định/kiểm tra theo cặp điểm đi--đến để tránh rò rỉ, và so sánh bằng kiểm định thống kê
    ghép cặp.
\end{enumerate}
